\section{The long-shift row and its geometric cap}\label{sec:long}

The row \(5W\) is the exact first-match remainder after \(4P\). Its source
keeps every earlier complement. In particular,
\[
 s\ge Y^{b_0},\quad K_{\rm res}\le Y^{\Delta_{\rm exp}},\quad
 \Delta_{\rm exp}=2\epsilon_0+k_0=3S_0.
\AutoTag{eq:old:7_1}
\]
The selected corners and both shifts keep their lower bounds from
Section~\ref{sec:partition}. The remainder predicate includes the exact
bound \(K_{\rm res}\le Y^{\Delta_{\rm exp}}\).

With \(e=\epsilon_0\), \(k=k_0\), \(l=l_0\), \(j=j_0\), and \(p=p_0\), write
\[
\begin{aligned}
 S_0=e=k&=\frac{30869}{5000000},&
 l&=2S_0,& j=p&=3S_0,\\
 b_0&=\frac{4166537}{30000000},&
 a_0&=\frac6{125},&
 \mu_{\rm HP}&=\frac{16}{25},\\
 \nu&=\frac1{64},&
 \xi&=\frac1{10000}.
\end{aligned}
\AutoTag{eq:old:7_2}
\]
Set
\[
 \begin{aligned}
 q_{\min}&=\frac23-S_0,\qquad q_{\max}=\frac23+2S_0,\\
 \tau_0&=1-\frac{3e}{2},\qquad
 \lambda=\log_Y(1+Y^{-1}),\qquad d_2=\log_Y2,\\
 \ell_C&=\log_Y(1+\sqrt3Y^{-\tau_0}),\qquad
 \ell_{\rm end}=\log_Y\bigl((1+\sqrt3Y^{-\tau_0})(1+2/Y)\bigr).
 \end{aligned}
\AutoTag{eq:old:7_3}
\]
The source inequalities imply \(q_{\min}<\log_Y M<q_{\max}+2\lambda\).
They also imply all endpoint comparisons used below. Each chart keeps its
own shift interval throughout the count.

\subsection{A polynomial on each ratio cell}

In a local chart box write
\[
 D_0\le J<2D_0,\quad E_0\le K<2E_0,\quad
 U_0\le f<2U_0,\quad S_h\le s<2S_h,\quad H_0=2S_h,
\]
\[
 X=D_0^2U_0,\qquad V_*=\frac{X}{E_0^2}.
\AutoTag{eq:old:7_4}
\]
These symbols are local box scales. On the actual population the source
identities give, uniformly,
\[
 \frac{X}{H_0}>\frac{Y^{\tau_0}}{96},\qquad
 \frac M{18}<U_0V_*\le4M,\qquad
 U_0\le f,\quad V_*\le4g.
\AutoTag{eq:old:7_5}
\]
The selected chart stays fixed as part of the first-match source label.
Independently, either counting orientation
\((J,K;f,g;s)=(A,D;ac,bd;u)\) or \((B,C;ad,bc;t)\) may be used
on that retained population. All constants in \eqref{eq:old:7_5} are
uniform in this orientation and its moving shift box.
Here \(T=fJ^2/s\) equals \(mn/(m+n+1)\ge m/2\) or
\(m(n+1)/(n-m)>m\), respectively. The dyadic definitions give
\(T<16X/H_0\). Hence \(X/H_0>m/32>Y^{\tau_0}/96\).
Also \(X/H_0<(Y+1)^4/Y^{b_0}<Y^4\) for sufficiently large \(Y\).
The conic gives the other comparisons in \eqref{eq:old:7_5} in
both orientations. It needs no lower bound for the unselected corners.

Put \(R_{\rm box}=D_0/E_0\), \(L_{\rm grid}=2N\), where
\[
 \bar m=\frac{q_{\max}}{1+3b_0^2}
       =\frac{203704280000000}{317360030572369},\qquad
 a_{\rm grid}=\bar m+\xi,\qquad N=\lceil Y^{a_{\rm grid}}\rceil.
\AutoTag{eq:old:7_6}
\]
The original grid consists of exactly the \(3N\) half-open cells
\[
 r_i=\frac{i}{2N}R_{\rm box},\qquad I_i=(r_i,r_{i+1}],
 \qquad i=N,\ldots,4N-1,
\AutoTag{eq:old:7_7}
\]
for the ratio \(J/K\). As \(D_0\le J<2D_0\) and
\(E_0\le K<2E_0\), these cells cover the whole local source ratio.
For \(r_0=r_i\), the relative width satisfies
\(r_{i+1}/r_i=1+1/i\le1+1/N\).

The exact determinant cap used to choose \(N\) is
\[
 \phi(\rho)=
 \begin{cases}
  9\rho/8,&0<\rho\le2/3,\\[2pt]
  \bigl(1+3(1-\rho)^2\bigr)^{-1},&2/3\le\rho<1.
 \end{cases}
\AutoTag{eq:old:7_8}
\]
Write \(q_M=\log_Y M\), \(c_0=\log_Y(D_0E_0)\), and
\(\tau=\log_Y(X/H_0)\). The limiting comparison ratio is
\(\rho=c_0/\tau\). The comparison
\(\log_YN>q_M\phi(\rho)+o(1)\) is the leading-order sufficient
condition. The proof below gives the uniform finite-\(Y\) margin after
the source and arithmetic collars. For the source comparison put
\(\sigma=\log_Ym\), \(z_s=\log_Ys\), and
\(\ell_L=\log_YL\). In the second chart also put
\(z_u=\log_Yu\). The exact identities from \eqref{eq:old:2_1}--\eqref{eq:old:2_2} used below are
\[
 AD={n-m\over t},\quad BC={m+n+1\over u},\quad
 L={(n-m)(m+n+1)\over Z},\quad Z^2M=m(m+1)n(n+1),
\]
\[
 A^2(ac)={umn\over m+n+1},\quad
 B^2(ad)={tm(n+1)\over n-m},\quad
 C^2(bc)={tn(m+1)\over n-m}.
\AutoTag{eq:old:7_8a}
\]
In the first chart, \(s=u\). The limiting source coordinates are
\[
 (c^*,\tau^*)=(\sigma+z_s-q_M/2,\sigma),
\qquad \rho_1^*={c^*\over\tau^*}
 =1+{z_s-q_M/2\over\sigma}.
\AutoTag{eq:old:7_8b}
\]
Before the change to local dyadic scales,
\[
 \log_Y(AD)-c^*
 =\log_Y{\sqrt{m(m+1)n(n+1)}\over m(m+n+1)},\quad
 \log_Y(A^2f/u)-\sigma=\log_Y{n\over m+n+1}.
\AutoTag{eq:old:7_8c}
\]
The first ratio factors as
\((1+1/m)^{1/2}[n/(m+n+1)](1+1/n)^{1/2}\). The middle
factor lies in \([1/(2+2/Y),1]\). The source gives
\(\tau_0-2d_2<\sigma\le1\). Both shifts satisfy
\(u,t\ge Y^{b_0}\), and
\(L=ut<\kappa_YY^{1/3+e}\). So for either chart shift \(s\in\{u,t\}\),
\[
 z_s\le q_{\max}/2-b_0+\lambda.
\AutoTag{eq:old:7_8b_common}
\]
On the selected chart, \(s=\min(u,t)\). So \(s^2\le L\) gives the stronger
individual estimate
\[
 z_s\le \frac16+\frac e2+\frac\lambda2.
\]
The alternative chart uses the other shift and the common bound
\eqref{eq:old:7_8b_common}.
The stronger estimate holds only for the selected smaller shift. It is
recorded for that source coordinate. It is not used in \eqref{eq:old:7_8e}.
Both chart caps use the common shift bound, which applies to either shift. The limiting
gap in this comparison is
\[
 \frac{q_{\max}}2-b_0-\left(\frac16+\frac e2\right)
 =\frac{30869}{1000000}>0.
\AutoTag{eq:old:7_8b_gap}
\]
The common shift bound minus the selected-shift bound is exactly
\[
 \frac{30869}{1000000}+\frac{\lambda}{2}>0.
\]
Thus their ordering needs no extra restriction on \(\lambda\). Small
\(\lambda\) is used below for the other strict endpoint comparisons. Its
terms tend to zero and are absorbed by their fixed margins.
Since \(q_M>q_{\min}\),
\[
 q_M/2-z_s\ge b_0-3e/2-\lambda>0
\]
for all sufficiently large \(Y\). Therefore
\[
 \frac{\partial\rho_1^*}{\partial\sigma}
 =\frac{q_M/2-z_s}{\sigma^2}>0,\qquad
 \frac{\partial\rho_1^*}{\partial z_s}=\frac1{\sigma}>0.
\]
Hence the maximum in \(\sigma\) is at \(\sigma=1\).
Fix \(z_s\). On the triangle branch
\[
 {d\over dq_M}\left[{9\over8}q_M(1-q_M/2+z_s)\right]
 ={9\over8}(1-q_M+z_s)>0.
\]
On the trapezoid branch, with \(r=q_M/2-z_s\),
\[
 {d\over dq_M}{q_M\over1+3(q_M/2-z_s)^2}
 ={1-3q_M^2/4+3z_s^2\over(1+3(q_M/2-z_s)^2)^2}>0,
\]
since \(q_M<1\). So this chart's maximum is at
\(q_M=q_{\max}+2\lambda\). At the joint upper endpoints
\(\sigma=1\), \(z_s=q_{\max}/2-b_0+\lambda\), the ratio is
\(\rho_1^*=1-b_0\). This is the alternative \(u\)-chart bound, even on a
source whose selected chart is \(t\). The comparison follows from the two
lower shift bounds and \(ut=L\).

In the second chart, \(s=t\). So the opposite shift satisfies
\(z_u\ge b_0\) on the entire \(5W\) source, and
\(\ell_L\le1/3+e+\lambda\). Its limiting coordinates are
\[
 (c^*,\tau^*)=(1-z_u,1-\ell_L+q_M/2),
 \qquad \rho_2^*={1-z_u\over1-\ell_L+q_M/2}.
\AutoTag{eq:old:7_8d}
\]
At \(\lambda=0\), increasing \(\ell_L\) and decreasing \(z_u\) increase
the ratio. So
\(\rho_2^*\le(1-b_0)/(q_{\min}+q_M/2)\). For finite \(Y\), the
denominator is at least \(D_\lambda=q_{\min}+q_M/2-\lambda\), where
\(D=q_{\min}+q_M/2\). With \(c=1-b_0\), the triangle-branch derivative
of \(q_M\phi(c/D_\lambda)\) is
\[
 \frac98\frac{c(q_{\min}-\lambda)}{D_\lambda^2}>0.
\]
Here \(9/8\) is the value of \(\phi(\rho)/\rho\) on this branch.
On the trapezoid branch the derivative has the sign of
\[
 1+3(1-\rho_2)^2-3(q_M/D_\lambda)\rho_2(1-\rho_2)
 >1-3\rho_2(1-\rho_2)\ge1/4,
\]
since \(q_M<D_\lambda\) for all sufficiently large \(Y\). The maximum
is at \(q_M=q_{\max}+2\lambda\). There \(D_\lambda=1\) and
\(\rho_2^*=1-b_0\). The same endpoint applies to the first chart above.
Thus both chart comparisons give the finite bound
\[
 q_M\phi(\rho_j^*)\le(q_{\max}+2\lambda)\phi(1-b_0)
 \le\bar m+2\lambda.
\AutoTag{eq:old:7_8e}
\]
Here \(j=1,2\) indexes the two local counting orientations and their
starred comparison ratios. Either orientation is available for every
source, whatever its selected-chart tag. The collar below compares the
local ratios with them.
Their common limiting cap as \(Y\to\infty\) is
\(q_{\max}/(1+3b_0^2)=\bar m\). They are alternative counting routes
for the same full \(5W\) source. Each source box is counted once, by its
better route. The choice is made separately for each box.

It remains to bound the finite collars. For
\(\rho\in(0,1)\), \(\phi'\le9/8\). On the triangle branch
\(\phi'=9/8\). On the trapezoid branch, with
\(r=1-\rho\in[0,1/3]\),
\(\phi'=6r/(1+3r^2)^2\le9/8\). Thus the partial derivatives of
\(q_M\phi(c_0/\tau)\) in \((q_M,c_0,\tau)\) are bounded by
\((1,9/4,9/4)\) whenever \(q_M<1\) and \(\tau>1/2\). In the two
charts, \eqref{eq:old:7_8a} and the local scale definitions give
\[
 |c_0-c^*|\le B_Y,\quad |\tau-\tau^*|\le B_Y,
 \qquad B_Y=6d_2+\tfrac12\ell_{\rm end}+\lambda.
\AutoTag{eq:old:7_8f}
\]
For chart one the exact residuals are \eqref{eq:old:7_8c}. For chart two
they are
\[
\begin{aligned}
 \log_Y(BC)-(1-z_u)&=\log_Y\frac{m+n+1}{Y},\\
 \tau^*-\log_Y(B^2f/t)
 &=\log_Y\left(\frac{Y}{m+n+1}
             \sqrt{\frac{m+1}{m}}\sqrt{\frac{n}{n+1}}\right).
\end{aligned}
\AutoTag{eq:old:7_8g}
\]
Here \(Y/2<m+n+1\le2Y\). So the logarithm of
\((m+n+1)/Y\), and also its negative, has magnitude at most \(d_2\).
In \eqref{eq:old:7_8c}, the middle factor \(n/(m+n+1)\) lies between
\(1/(2+2/Y)\) and \(1\). Its logarithm has magnitude at most
\(d_2+\lambda\). The other factors contribute at most
\((h_m+h_n)/2\). The local dyadic replacement of \(AD\) by \(D_0E_0\)
costs at most \(2d_2\). The replacement of \(A^2f/u\) by \(X/H_0\)
costs at most \(4d_2\). In the second chart, the first residual in
\eqref{eq:old:7_8g} costs at most \(d_2\), and the second at most
\(d_2+(h_m+h_n)/2\). Its local scale changes have the same respective
costs \(2d_2\) and \(4d_2\). Also \(h_m+h_n\le\ell_{\rm end}\), where
\(h_m=\log_Y(1+1/m)\), \(h_n=\log_Y(1+1/n)\). The bounds
\(m>Y^{\tau_0}/\sqrt3\) and \(n>Y/2\) imply
\[
 h_m+h_n\le
 \log_Y\bigl((1+\sqrt3Y^{-\tau_0})(1+2/Y)\bigr)
 =\ell_{\rm end}.
\]
The larger of these bounds gives \eqref{eq:old:7_8f}, with room for the
extra \(\lambda\) term in the first chart. This keeps one
\(d_2\) unit of slack beyond the largest bound above.

Put
\[
 E_Y=\rho_Y+\log_Y(96\cdot18\cdot324\cdot8\cdot3\cdot2),
 \qquad \rho_Y=11d_2+2\ell_{\rm end}+2\lambda.
\AutoTag{eq:old:7_8h}
\]
The six factors pay for fixed comparisons. The factor 96 is used
in \(X/H_0>Y^{\tau_0}/96\). The factor 18 is used in \(M<18U_0V_*\).
The factor 324 is used in
the first anchor \(\Lambda_1=\log(324U_0V_*)\), together with
\(U_0V_*\le4M\). The factor 8 is used in the second anchor
\(\Lambda_2=\log(8D_0E_0)\). The factors 3 and 2 are used in \(X_1=N/3\)
and in \(X_2=X/(2H_0)\). The source collars contribute at most
\(2\lambda\) in \(q_M\). Then \eqref{eq:old:7_8f} contributes at most
\((1+9/4+9/4)B_Y+2\lambda\le6\rho_Y\). With the fixed logarithmic
comparisons above, these terms are bounded by \(6E_Y\).
The upper comparison \(U_0V_*\le4M\) contributes \(\log_Y4=2d_2\).
This lies in the \(11d_2\) part of \(\rho_Y\), hence in the
\(6\rho_Y\) collar. Also
\[
 \frac{\Lambda_1}{\log Y}=q_M+\log_Y(324U_0V_*/M),
\]
so the coefficient of \(q_M=\log_YM\) in this anchor is exactly one.
Hence the full finite cap is at most \(\bar m+6E_Y\). Since \(E_Y\to0\),
we may impose \(6E_Y<\xi/2\). This leaves a grid margin of at least
\(\xi/2\) in \eqref{eq:old:7_6}.

On the trapezoid branch the strict determinant surplus
at the limiting source comparison is at least
\[
\begin{aligned}
 \frac{q_{\min}\xi}{\bar m(\bar m+\xi)}
 &=\frac{997846754805373553496571768583986273}
         {6225284767495687836225589800000000000000}\\
 &=0.0001602893\ldots>0.
\end{aligned}
\AutoTag{eq:old:7_9}
\]
The triangle branch has its own strict surplus. The fixed margin \(\xi\)
is charged in \eqref{eq:old:7_6}, before the degree family is chosen. The
closed source-parameter ranges are compact, and \eqref{eq:old:7_9} gives a
positive surplus. Together they give a fixed \(\delta_T>0\) below the
minimum trapezoid-branch gap. The inequalities
\[
 q_{\min}-a_{\rm grid}>0,\quad \tau_0-a_{\rm grid}>0,\quad
 \tau_0-q_{\max}=\frac{9351751}{30000000}>0,\quad
 b_0-\frac{3e}{2}=\frac{972179}{7500000}>0
\AutoTag{eq:old:7_10}
\]
put all ratios in compact subintervals of their required ranges. Thus one
finite set of degrees, one Taylor order, and one threshold work for every
cell and shift.

\begin{lemma}[Ratio-cell determinant]\label{lem:ratio-cell-determinant}
Fix the parameters in \eqref{eq:old:7_2}--\eqref{eq:old:7_6}, with
\(3/5<a_{\rm grid}<\min(q_{\min},\tau_0)\),
\(q_{\max}<\tau_0\), and \(b_0>3e/2\). Let \(Y\) be sufficiently large.
Keep the exact first-match \(5W\) population
\eqref{eq:old:3_1}, with its source masks, original selected-chart labels,
and all preceding-row complements. Take either counting orientation,
any local box satisfying
\eqref{eq:old:7_4}--\eqref{eq:old:7_5} and \(X/H_0\le Y^4\), and any
cell in \eqref{eq:old:7_7}. Then there is a
nonzero polynomial
\(F_i\in\mathbb Z[X,Z]\) of degree at
most \(D_\theta\) and coefficient height at most \(Y^{H_\theta}\). Here
\(D_\theta,H_\theta\) depend only on the fixed parameter vector. It
satisfies \(F_i(J/K,g/f)=0\) for every retained source point in that cell
and every admissible shift in the box. The bounds are uniform in the
counting orientation, original selected chart, box, cell, labels and shift.
\end{lemma}
\begin{proof}
The proof is an integer calculation. We give the form used below. Put
\[
 x={f\over g},\quad y={K\over J},\quad y_0=r_i^{-1},\quad
 \delta_{\rm cell}={s\over fJ^2},\quad
 \alpha_{\rm cell}={x\over y_0^2}-1.
\]
The conic gives \(x=y^2/(1+\delta_{\rm cell})\) and
\(y/y_0=\sqrt{1+\alpha_{\rm cell}}\sqrt{1+\delta_{\rm cell}}\).
As \(J/K\in(r_i,r_{i+1}]\), we have
\(i/(i+1)\le y/y_0<1\), with \(i\ge N\). Also
\[
0<\delta_{\rm cell}\le \frac{H_0}{U_0D_0^2}
=\frac{H_0}{X}\le\frac1N
\]
for sufficiently large \(Y\), since \(X/H_0>Y^{\tau_0}/96\),
\(N\le2Y^{a_{\rm grid}}\), and \(a_{\rm grid}<\tau_0\). Therefore
\[
|\alpha_{\rm cell}|
=1-\frac{(y/y_0)^2}{1+\delta_{\rm cell}}
\le \frac2N+\frac{H_0}{X}
\le\frac3N.
\]
Next,
\[
0\le\frac{y}{y_0}-\sqrt{1+\alpha_{\rm cell}}
=\sqrt{1+\alpha_{\rm cell}}\bigl(\sqrt{1+\delta_{\rm cell}}-1\bigr)
\le\frac{H_0}{2X},
\]
using \(\sqrt{1+\alpha_{\rm cell}}\le1\) and
\(\sqrt{1+\delta}-1\le\delta/2\). Thus
\[
 |\alpha_{\rm cell}|\le3/N,\qquad
 0\le y/y_0-\sqrt{1+\alpha_{\rm cell}}\le H_0/(2X).
 \AutoTag{eq:old:7_10a}
\]
As \(X_2=X/(2H_0)\), the conic error is
\(H_0/(2X)=X_2^{-1}/4\).
Set \(T_{\rm tay}=7\). Let \(P_{T_{\rm tay}}\) be the degree-seven
Taylor polynomial of \(\sqrt{1+\alpha_{\rm cell}}-1\). For
\(c_k=|\binom{1/2}{k}|\), we have
\(c_{k+1}/c_k=(k-1/2)/(k+1)<1\) for \(k\ge1\), and
\(c_8=429/32768\). So, for
\(|\alpha_{\rm cell}|\le1/2\), the remainder is at most
\[
 \frac{c_8|\alpha_{\rm cell}|^8}{1-|\alpha_{\rm cell}|}
 \le\frac{429}{16384}|\alpha_{\rm cell}|^8
 <2|\alpha_{\rm cell}|^8.
\]
The source gives
\(X/H_0\le Y^4\) and \(a_{\rm grid}>3/5\). Hence, once
\(4\cdot3^8Y^{-8a_{\rm grid}}\le Y^{-4}\), the Taylor remainder is at
most \(X_2^{-1}/2\), where \(X_1=N/3\), \(X_2=X/(2H_0)\). To check this,
\(X/H_0\le Y^4\) gives \(X_2^{-1}\ge2Y^{-4}\). So the
Taylor bound is at most \(Y^{-4}/2\le X_2^{-1}/2\). Add the
conic error \(X_2^{-1}/4\). Then
\(\beta_{\rm cell}=y/y_0-1-P_{T_{\rm tay}}(\alpha_{\rm cell})\) satisfies
\(|\alpha_{\rm cell}|\le X_1^{-1}\),
\(|\beta_{\rm cell}|\le X_2^{-1}\), uniformly.

Take columns \(x^r y^j\), \(0\le r\le A_{\rm det}\),
\(0\le j\le B_{\rm det}\). Multiply the row for \((f,g,J,K)\) by
\(g^{A_{\rm det}}J^{B_{\rm det}}\). The entries are the integers
\[
 f^r g^{A_{\rm det}-r}K^jJ^{B_{\rm det}-j}.             \AutoTag{eq:old:7_10b}
\]
After the anchor is removed the entries are
\((1+\alpha_{\rm cell})^r
(1+P_{T_{\rm tay}}(\alpha_{\rm cell})+\beta_{\rm cell})^j\). With
\(D_c=(A_{\rm det}+1)(B_{\rm det}+1)\), \(g<9V_*\),
\(y_0\le2E_0/D_0\), and \(J<2D_0\),
the total row and column arithmetic cost is at most
\[
 {D_c\over2}\{A_{\rm det}\Lambda_1+B_{\rm det}\Lambda_2\},\qquad
 \Lambda_1=\log(324U_0V_*),\quad
 \Lambda_2=\log(8D_0E_0).                                \AutoTag{eq:old:7_10c}
\]
This includes the factor
\(y_0^{D_c(A_{\rm det}+B_{\rm det}/2)}\).
To verify \(g<9V_*\), multiply the conic
\(g/f=(J/K)^2+s/(fK^2)\) by \(f\):
\[
g=f(J/K)^2+s/K^2
 <8U_0D_0^2/E_0^2+H_0/E_0^2
 =8V_*+H_0/E_0^2<9V_*.
\]
Here \(f<2U_0\), \(J/K<2D_0/E_0\), \(s<H_0\), \(K\ge E_0\),
and \(X/H_0>Y^{\tau_0}/96>1\) for sufficiently large \(Y\).

For the finite Cauchy--Binet calculation let
\(\ell_1=\log X_1\), \(\ell_2=\log X_2\),
\(p_1=\Lambda_1/\ell_1\), \(p_2=\Lambda_2/\ell_2\), and
\(\ell_{\rm det}=\ell_1/\ell_2\). Set \(\ell=\ell_{\rm det}\).
Here \(X_1=N/3\), \(X_2=X/(2H_0)\),
\(\ell_1=\log(N/3)\), and \(\ell_2=\log(X/(2H_0))\). So the finite
parameters in the strip comparison are exactly
\[
 p_1=\frac{\log(324U_0V_*)}{\log(N/3)},\qquad
 p_2=\frac{\log(8D_0E_0)}{\log(X/(2H_0))},\qquad
 \ell_{\rm det}=\frac{\log(N/3)}{\log(X/(2H_0))}.
\]
Define the strip weight
\[
 W_{\rm strip}=\min_{\substack{\mathcal S\subset\mathbb Z_{\ge0}\times\{0,\ldots,B_{\rm det}\}\\
                                  |\mathcal S|=D_c}}
                  \sum_{(i,k)\in\mathcal S}(i\ell_1+k\ell_2).
\]
After column normalization the entries are polynomials in
\((\alpha_{\rm cell},\beta_{\rm cell})\) of \(\beta_{\rm cell}\)-degree
at most \(B_{\rm det}\). The Taylor coefficients of
\(\sqrt{1+\alpha_{\rm cell}}-1\) have absolute sum at most one. So each
column's coefficient \(\ell^1\)-norm is at most \(2^{A_{\rm det}}3^{B_{\rm det}}\).
Cauchy--Binet, Hadamard, and the at most
\(2^{(A_{\rm det}+1+B_{\rm det}T_{\rm tay})(B_{\rm det}+1)}\) subsets
of the finite monomial set
then give
\[
 \begin{aligned}
 \log|\Delta_1|&\le \log K_{A_{\rm det},B_{\rm det},T_{\rm tay}}-W_{\rm strip},\\
 K_{A_{\rm det},B_{\rm det},T_{\rm tay}}
 &=2^{(A_{\rm det}+1+B_{\rm det}T_{\rm tay})(B_{\rm det}+1)}(D_c!)^2
 (2^{A_{\rm det}}3^{B_{\rm det}})^{D_c}.
 \end{aligned}
\AutoTag{eq:old:7_10d}
\]
Here \(K_{A_{\rm det},B_{\rm det},T_{\rm tay}}\) is fixed once the degree
and Taylor order are fixed.
It is independent of \(Y\), the box, the cell, and the shift.

We give the strip lower bound explicitly. Put
\(A'=A_{\rm det}+1\), \(B'=B_{\rm det}+1\), so \(D_c=A'B'\).
Suppose
\(A'\ell_1\ge B'\ell_2/2\). To each lattice point
\((i,k)\) attach its unit square \([i,i+1)\times[k,k+1)\). For
\(\Phi(x,y)=x\ell_1+y\ell_2\), the integral over that square is
\(i\ell_1+k\ell_2+(\ell_1+\ell_2)/2\). The union \(E\) of the
\(D_c\) squares has area \(D_c\) and lies in the strip
\(x\ge0,0\le y\le B'\). Among all sets of this area in that strip, the
integral of \(\Phi\) is smallest on the sublevel set
\(E_* = \{\Phi\le w\}\cap\{x\ge0,\ 0\le y\le B'\}\). Points outside \(E_*\)
have \(\Phi\ge w\), and points removed from \(E_*\) have \(\Phi\le w\).
So exchanging equal areas cannot lower the integral. The condition above
makes \(E_*\) a trapezoid with right edge
\(x=(w-y\ell_2)/\ell_1\). Its area is
\[
 \frac{wB'-\ell_2B'^2/2}{\ell_1}=A'B',\qquad
 w=A'\ell_1+B'\ell_2/2.
\]
Integrating first in \(x\), then in \(y\), gives
\[
 \int_{E_*}\Phi
 =\int_0^{B'}\frac{w^2-y^2\ell_2^2}{2\ell_1}\,dy
 =\frac12A'^2B'\ell_1+\frac12A'B'^2\ell_2
   -\frac{B'^3\ell_2^2}{24\ell_1}.
\AutoTag{eq:old:7_10e}
\]
Subtracting the unit-square correction gives
\[
 W_{\rm strip}\ge \frac12A'^2B'\ell_1+
 \frac12A'B'^2\ell_2-\frac{B'^3\ell_2^2}{24\ell_1}
 -\frac12A'B'(\ell_1+\ell_2).
\AutoTag{eq:old:7_10f}
\]

The integer row clearing and the anchor factors have total logarithmic
cost at most
\(C_{\rm arith}=D_c(A_{\rm det}\Lambda_1+B_{\rm det}\Lambda_2)/2\), as in
\eqref{eq:old:7_10c}. As \(A_{\rm det}\le A'\) and \(B_{\rm det}\le B'\),
this is at most
\(D_c(A'\Lambda_1+B'\Lambda_2)/2\). Put \(z=A'/B'\). Division by
\(B'^3\ell_2/2\) gives, on the trapezoid branch,
\[
 W_{\rm strip}-C_{\rm arith}\ge\frac{B'^3\ell_2}{2}
 \left[-\ell_{\rm det}(p_1-1)z^2+(1-p_2)z-\frac1{12\ell_{\rm det}}
       -\frac{z(\ell_{\rm det}+1)}{B'}\right].          \AutoTag{eq:old:7_10g}
\]
The first three terms form the concave quadratic
\[
 G(z)=-\ell_{\rm det}(p_1-1)z^2+(1-p_2)z
       -\frac1{12\ell_{\rm det}}.
\]
If \(p_2\ge2/3\), \(p_1>1\), and
\(p_1-1\le3(1-p_2)^2-\delta_T\), then
\[
 z_* =\frac{1-p_2}{2\ell_{\rm det}(p_1-1)}
 \ge\frac1{2\ell_{\rm det}},\qquad
 G(z_*)=\frac{3(1-p_2)^2-(p_1-1)}
 {12\ell_{\rm det}(p_1-1)}>0.
\]
Here we use \(3(1-p_2)^2\le1-p_2\) for \(p_2\ge2/3\). On the compact
source range, \(z_*\), \(\ell\), and \(p_1-1\) are uniformly bounded,
and \(\ell(p_1-1)\) is bounded away from zero. Choose one fixed \(B'\)
so large that the rounding loss
\(\ell(p_1-1)/B'^2\) and the boundary term
\((z_*+1/B')(\ell+1)/B'\) sum to less than half the minimum of
\(G(z_*)\). Then set \(A'=\lceil z_*B'\rceil\). Since
\(G(z)=G(z_*)-\ell(p_1-1)(z-z_*)^2\), the bracket in \eqref{eq:old:7_10g} is
positive uniformly. The strict source margin in \eqref{eq:old:7_9} gives the
positive minimum. Choose the fixed lower threshold so that the total
\(O(1/\log Y)\) collar is at most half that minimum. The finite comparison
then keeps at least half the limiting gap uniformly.

For \(p_2\le2/3\), use the whole quadrant in place of the strip.
The sublevel set \(\Phi\le w\) in \(x,y\ge0\) is a triangle of area
\[
 \frac{w^2}{2\ell_1\ell_2}=D_c,\qquad
 \int_{E_*}\Phi=\frac{2\sqrt2}{3}D_c^{3/2}\sqrt{\ell_1\ell_2}.
\]
The unit-square error is \(D_c(\ell_1+\ell_2)/2\). Put \(z=A'/B'\). The
leading weight minus the upper bound for \(C_{\rm arith}\), divided by
\(B'^3\ell_2/2\), is
\[
 z\left(\frac{4\sqrt2}{3}\sqrt{\ell z}-p_1\ell z-p_2\right).
\]
Set \(v=\sqrt{\ell z}\). The parenthesis is the concave quadratic
\((4\sqrt2/3)v-p_1v^2-p_2\). Its maximum is at
\(v=2\sqrt2/(3p_1)\), which corresponds to
\[
 z_* =\frac{8}{9p_1^2\ell},\qquad
 A'=\left\lceil\frac{8B'}{9p_1^2\ell}\right\rceil,
\]
and the maximum value is
\((8/9-p_1p_2)/p_1\). Thus the triangle branch has a uniform positive
gap whenever \(p_1p_2<8/9\). A fixed large enough \(B'\) absorbs
the square and rounding errors. At \(p_2=2/3\), the two threshold
conditions agree. In both branches the source inequality
\(N>Y^{q\phi(\rho)}\), with the fixed margin \eqref{eq:old:7_9}, gives
these strict conditions for every cell.

After the fixed collar cutoff, the normalized bracket has a uniform
positive lower bound on both branches. Also
\(\ell_2=\log(X/(2H_0))>\tau_0\log Y-\log192\). Hence
\(W_{\rm strip}-C_{\rm arith}\ge c_\theta\log Y\) for some fixed
\(c_\theta>0\) and all sufficiently large \(Y\). The degrees and
Taylor order range over a fixed finite family. So
\[
 K_{\max}=\max K_{A_{\rm det},B_{\rm det},7}
\]
is fixed. Increase the determinant cutoff until
\(c_\theta\log Y>\log K_{\max}\). Then
\(W_{\rm strip}-C_{\rm arith}>\log K_{\max}\) uniformly. Every
cleared \(D_c\)-minor is an integer of absolute value less than one,
hence zero. The matrix has rank \(r<D_c\). If \(r=0\),
a coordinate vector gives a kernel polynomial. Otherwise choose \(r\)
pivot columns and \(r\) rows giving a nonsingular minor \(A\), then one
free column \(j\). Set the \(j\)-coordinate to \(\det A\), the other
free coordinates to zero, and the pivot coordinates to
\(-\operatorname{adj}(A)M_{I,j}\). The resulting nonzero integer vector
annihilates the selected rows. They span the row space, so it annihilates
every row. Its coordinates are bounded-order minors, so they have
polynomial height. This gives a nonzero integer polynomial \(C(x,y)\)
vanishing at
\((x,y)=(f/g,K/J)\). If a cell has fewer than \(D_c\) rows, then
\(r<D_c\) is automatic and its points are counted directly; this includes
an empty cell. Define
\[
 F_i(R,S)=R^{B_{\rm det}}S^{A_{\rm det}}C(1/S,1/R).
\AutoTag{eq:old:7_10h}
\]
This is a nonzero integer polynomial of degrees at most
\((B_{\rm det},A_{\rm det})\), with
the same coefficient height. It vanishes at
\((R,S)=(J/K,g/f)\). The substitution permutes the coefficients, so the
integral polynomial and its height are unchanged. This proves the
ratio-cell lemma.
\end{proof}

\subsection{One reduced lattice and the complete cell census}

The auxiliary polynomial is bihomogenized in the two primitive pairs.
For its high-degree factors use the single canonical coefficient lattice
\[
 j(i)=\left\lfloor\frac{i^2}{L_{\rm grid}}\right\rfloor,\quad
 t_i=\frac{j(i)}{L_{\rm grid}}\frac{V_*}{U_0},\quad
 W_i(g,f)=\left(\frac{L_{\rm grid}}{V_*}(g-t_if),\frac f{U_0}\right).
\AutoTag{eq:old:7_11}
\]
Its determinant is \(L_{\rm grid}/(U_0V_*)\). The map \(i\mapsto j(i)\)
is injective, because
\(((i+1)^2-i^2)/L_{\rm grid}=(2i+1)/(2N)>1\).
Recall \(R_{\rm box}=D_0/E_0\). So \(V_*/U_0=R_{\rm box}^2\), and
\[
 W_i=\begin{pmatrix}L_{\rm grid}/V_*&-j(i)/U_0\\0&1/U_0\end{pmatrix}.
\]
The source conic gives
\(g/f=(J/K)^2+s/(fK^2)\). Since
\(r_i=iR_{\rm box}/L_{\rm grid}\), \(J/K\le r_i(1+1/i)\), and
\(0\le r_i^2-t_i<R_{\rm box}^2/L_{\rm grid}\), for \(N\ge4\) we get
\[
 0<g/f-t_i\le R_{\rm box}^2\left({10\over N}+{H_0\over X}\right).
\]
As \(f<2U_0\), \(L_{\rm grid}=2N\), and
\(NH_0/X\le2^{-20}\), the first coordinate of \(W_i(g,f)\) lies
between zero and \(41\). Its second lies in \([1,2]\). So all
source images lie in one fixed bounded rectangle. Put
\(\Lambda_i=W_i\mathbb Z^2\). Choose a Lagrange-reduced basis
\(v_{1,i},v_{2,i}\) of this image lattice, with lengths
\(\lambda_{1,i}\le\lambda_{2,i}\) and angle bounded away from zero.
Write \(v_{j,i}=W_i h_{j,i}\). Since \(v_{1,i},v_{2,i}\) are a basis
of \(\Lambda_i\), their integer preimages \(h_{1,i},h_{2,i}\) form a
unimodular basis of \(\mathbb Z^2\). The vector \(v_{1,i}\) is shortest
in the Euclidean norm of the image lattice \(\Lambda_i\). Define
\[
 T_{1,i}={32\sqrt2\over\lambda_{1,i}},\qquad
 T_{2,i}={32\sqrt2\over\lambda_{2,i}}.
\]
Then \(T_{2,i}\le T_{1,i}\). The determinant identity gives
\(T_{1,i}T_{2,i}\asymp U_0V_*/L_{\rm grid}\). A unimodular change
preserves primitivity. By the bounded rectangle and the reduced-basis
angle, every source point has coordinates
\(O(T_{1,i})\) in this basis. The map and its inverse are rational matrices of
polynomial height. Explicitly,
\[
 W_i^{-1}(y_1,y_2)=
 \left({V_*\over L_{\rm grid}}(y_1+j(i)y_2),U_0y_2\right).
\]
A common denominator for \(W_i\) is at most a fixed power of \(Y\). So
every nonzero vector of \(\Lambda_i\) has length at least \(Y^{-C}\).
Its covolume is \(L_{\rm grid}/(U_0V_*)\le Y^C\). Minkowski's bound
and the reduced-angle determinant identity then give
\(\lambda_{1,i},\lambda_{2,i}\le Y^C\). The inverse matrix gives
\(\lVert h_{j,i}\rVert\le Y^{C'}\). Thus the unimodular basis and every
fixed-degree transformed factor have polynomial height, uniformly in the
box, cell and labels. The first-coordinate radius \(T_{1,i}\) satisfies
\[
 \sum_i(T_{1,i}+1)\ll_\eta
 Y^\eta\{N+\sqrt{U_0V_*N}+U_0+V_*\}.
\AutoTag{eq:old:7_12}
\]
Take a dyadic range \(D_{\rm lat}<T_{1,i}\le
2D_{\rm lat}\). Write the primitive shortest preimage in source
coordinates as \((\omega_1,\omega_2)\). It satisfies
\[
 |\omega_1|\ll V_*/D_{\rm lat},\quad |\omega_2|\ll U_0/D_{\rm lat},\quad
 \omega_2j(i)=\frac{L_{\rm grid}U_0}{V_*}\omega_1+O(U_0/D_{\rm lat}).
\AutoTag{eq:old:7_13}
\]
The first coordinate bound follows from the first component of \(W_i\),
the bound \(j(i)/L_{\rm grid}=O(1)\), and
\(|\omega_2|\ll U_0/D_{\rm lat}\). Suppose both coordinates are nonzero.
Fix \(\omega_1\) first. There are \(O(V_*/D_{\rm lat})\) choices. For each such
\(\omega_1\), \eqref{eq:old:7_13} places the integer product
\(w_i=\omega_2j(i)\) in an interval
of length \(O(U_0/D_{\rm lat})\). As \(\omega_2\ne0\),
\(U_0/D_{\rm lat}\gg1\). So this interval contains
\(O(U_0/D_{\rm lat})\) integers. Also \(w_i\ne0\), because
\(j(i)>0\).
For each fixed \(w_i\), the divisor bound gives \(Y^\eta\) choices of
\((\omega_2,j(i))\). Injectivity of \(j\) gives at most one index \(i\) for
each \(j(i)\). So there are
\(O_\eta(Y^\eta U_0V_*/D_{\rm lat}^2)\) indices in this dyadic range.
Weighting by \(D_{\rm lat}\) and summing from the lattice minimum gives
\(O_\eta(Y^\eta\sqrt{U_0V_*N})\). If \(\omega_2=0\),
primitivity gives
\(\omega_1=\pm1\), and the total is \(O(V_*+N)\). If \(\omega_1=0\), it is
\(O(U_0+N)\). The cells with \(T_{1,i}<1\) and all endpoint \(+1\)
terms are included in \(N\). This proves \eqref{eq:old:7_12}.

The nonzero product is at most a fixed power of \(Y\).
So its divisor estimate is uniform over the boxes and cells. The lower
limit \(D_{\rm lat}\gg\sqrt{U_0V_*/N}\) follows from
\(T_{1,i}T_{2,i}\asymp U_0V_*/N\) and \(T_{1,i}\ge T_{2,i}\).
Summing the unweighted number of non-axis indices over the same dyadic
ranges gives \(O_\eta(Y^\eta N)\). On the first axis, \(\omega_2=0\).
There primitivity gives \(\omega_1=\pm1\), and the contribution is
\(V_*+N\). On
the second axis, \(\omega_1=0\). There it gives \(\omega_2=\pm1\), and the
contribution is \(U_0+N\). The \(T_{1,i}<1\) bins contribute another
\(O(N)\). The bound \eqref{eq:old:7_12} keeps all these terms.
The bins with \(T_{2,i}<1\) are included in the \(+1\) terms.

A factor has degree \((a_c,b_q)\): corner degree first,
coefficient degree second. If \(a_c\ge2\) and \(b_q\ge1\), apply
Heath--Brown's published Lemma~2. Put the primitive coefficient-lattice
pair first and the primitive corner pair second. The bound is
\[
 O_{a_c,b_q,\eta}\bigl((T_{1,i}+1)^{2/a_c+\eta}
 \|F\|^\eta\bigr),\qquad 2/a_c\le1,
\AutoTag{eq:old:7_14}
\]
and the second pair is unrestricted. The statement in
Appendix~\ref{app:sources}, cited as
\cite[Lemma~2, p.~8]{HB2012}, applies with the reduced coefficient-lattice
pair first and the corner pair second. Reversing the degree order gives
second-pair degree \(a_c\). The coefficient pair is bounded by
\(X\asymp T_{1,i}+1\), while the corner pair is unrestricted. The source
gcd conditions make both pairs primitive. The factor height is
\(Y^{O(1)}\) before \(\eta\) is chosen. The stated bound therefore gives
\eqref{eq:old:7_14}. Polynomial factor heights and the finite degree family
let the loss be made arbitrary.

The remaining factors are as follows.
\begin{enumerate}
\item A rational factor that is not absolutely irreducible has rational
points only on intersections with distinct conjugate components.
Resultants in the four affine and infinity charts give \(O(1)\) points at
fixed degree.
\item If \(a_c=0\), a coefficient pair fixes the polynomial but not yet
the shift. First choose its ordered residual factors. These fix the
nonzero integer \(L=cd-ab\). Then enumerate \(s\mid L\). Apply the
uniform square-fibre theorem of Appendix~\ref{app:fibre} with
\((a,b,h,n,m)=(f,g,s,J,K)\), \(x=2Y\), and fixed height exponent
\(A_{\rm ht}=2\). This costs \(Y^\eta\), including the shift divisor count.
\item If \(b_q=0\), fix \((J,K)\) and compare two coefficient pairs:
\[
\begin{aligned}
 K^2(g_1f_2-g_2f_1)&=s_1f_2-s_2f_1,\\
 |g_1f_2-g_2f_1|
 &<\frac{3S_hU_0}{E_0^2}
 =\frac{3U_0V_*H_0}{2X}
 \le\frac{6MH_0}{X}\\
 &<576Y^{q_{\max}+2\lambda-\tau_0}=o(1).
\end{aligned}
\AutoTag{eq:old:7_15}
\]
The determinant is an integer, so it vanishes for all sufficiently large
\(Y\). Primitivity forces the coefficient pairs to agree. Then the conic
fixes the shift. The strict gap is
\(\tau_0-q_{\max}=9351751/30000000\). Since
\(s_if_j\in[S_hU_0,4S_hU_0)\), the difference of two such products
has absolute value less than \(3S_hU_0\). Then \eqref{eq:old:7_5} gives
the bound.
\item For \(b_q\ge1\) and \(a_c=1\), write the factor as
\(K\mathcal P(f,g)+J\mathcal Q(f,g)\), with coprime homogeneous
\(\mathcal P,\mathcal Q\). Their resultant is nonzero. By its Bezout
identities,
\(\gcd(\mathcal P(f,g),\mathcal Q(f,g))\) divides a fixed
polynomial-height integer. Write
\(d_0=\gcd(\mathcal P(f,g),\mathcal Q(f,g))\). Homogeneous Bezout
identities and \(\gcd(f,g)=1\) give
\(d_0\mid\operatorname{Res}(\mathcal P,\mathcal Q)\). Primitivity
then determines the positive corner pair
\((J,K)=\pm(\mathcal P(f,g)/d_0,-\mathcal Q(f,g)/d_0)\).

Let \(d_f=f_J\) be the full prime-power part of \(f\) supported on primes
dividing \(J\). Let \(d_g=g_K\) be the full prime-power part of \(g\)
supported on primes dividing \(K\). Put \(v=f/d_f\), \(w=g/d_g\).
As \(d_f,d_g\) contain the full powers of their
supported primes, \(d_f,d_g,v,w\) are pairwise coprime, and
\(d_fd_gvw=fg=M\).
The four endpoint patterns use the labels
\[
\begin{array}{c|c|c|c}
 \mathcal P(0,1)&\mathcal Q(1,0)&\ell_f&\ell_g\\ \hline
 \ne0&\ne0&\rad(d_f)&\rad(d_g)\\
 \ne0&=0&\rad(d_f)&w\\
 =0&\ne0&v&\rad(d_g)\\
 =0&=0&v&w .
\end{array}                                                   \AutoTag{eq:old:7_16}
\]
Suppose \(\mathcal P(0,1)\ne0\). Every prime dividing \(d_f\) divides \(J\)
and does not divide \(K\) or \(g\). The factor equation modulo that
prime gives \(\mathcal P(f,g)\equiv0\). Homogeneity then gives
\(\mathcal P(0,1)\equiv0\). Hence \(\rad(d_f)\mid\mathcal P(0,1)\).
If \(\mathcal P(0,1)=0\), then \(f\mid\mathcal P(f,g)\). As
\(\gcd(v,J)=1\), the factor equation gives \(v\mid\mathcal Q(f,g)\).
So \(v\mid d_0\mid\operatorname{Res}(\mathcal P,\mathcal Q)\).
The same argument with \(g,K,w,\mathcal Q(1,0)\) proves the two
\(g\)-side entries. These divisibilities include full prime powers in
\(v,w\). The zero endpoint cases are handled by the matching table
entries. So each label divides a fixed nonzero endpoint coefficient or
resultant of polynomial height. Hence there are \(\ll_\eta Y^\eta\)
possible pairs of endpoint labels for any fixed \(\eta>0\).

For each of the four patterns, the sum shown is the lower-corner
base-\(Y\) exponent of the matching full assigned-part product. In
chart 1 these sums are
\[
 P+Q,\quad P+(q-c_1-Q),\quad(c_1-P)+Q,\quad
 (c_1-P)+(q-c_1-Q),
\]
and in chart 2 they are
\[
 R_{\rm LP}+T,\quad R_{\rm LP}+(q-c_2-T),\quad
 (c_2-R_{\rm LP})+T,\quad
 (c_2-R_{\rm LP})+(q-c_2-T),
\]
in the same row order. Here \(P,Q,R_{\rm LP},T\) are the logarithmic sums
in \eqref{eq:old:7_18}. The forms are \(\mathcal P,\mathcal Q\). Put
\(k(x)=x/\rad(x)\). Pairwise coprimality gives
\(k(d_f)k(d_g)k(v)k(w)=M/\rad(M)=K_{\rm res}\). A label is either the
full assigned part or its radical. So in every row
\[
 \ell_f\ell_g\ge
 \frac{\text{corresponding assigned-part product}}{K_{\rm res}}.
\]
By \(K_{\rm res}\le Y^{\Delta_{\rm exp}}\), each matching actual
product is at least \(Y^{E_{\rm pat}}\), where \(E_{\rm pat}\) is its
lower-corner exponent. Thus
\(\ell_f\ell_g\ge Y^{E_{\rm pat}-\Delta_{\rm exp}}\). These are exactly
the four chart pattern terms with loss \(-\Delta_{\rm exp}\) in
\eqref{eq:old:7_19}.
Each of the eight assigned factors is below twice its dyadic lower
endpoint. So \(M\le2^8Y^q\). Also \(U_0V_*\le4M\). Hence
\[
 U_0V_*\le2^{10}Y^q,\qquad
 \frac{U_0V_*}{\ell_f\ell_g}
 \le2^{10}Y^{q-E_{\rm pat}+\Delta_{\rm exp}}.
\]
For the spacing estimate, the conic gives
\[
 {g\over f}=\left({J\over K}\right)^2+{s\over fK^2}.
\]
Over a cell of width \(O(D_0/(E_0N))\) in \(J/K\), the first term
varies by \(O(D_0^2/(E_0^2N))=O(V_*/(U_0N))\). The second term has
size \(O(H_0/(U_0E_0^2))\), smaller by the factor
\[
 {H_0/(U_0E_0^2)\over V_*/(U_0N)}={NH_0\over X}=o(1)
\]
by \eqref{eq:old:7_5}--\eqref{eq:old:7_6} and \(a_{\rm grid}<\tau_0\).
Hence the ratios \(g/f\)
in one cell lie in an interval of length \(O(V_*/(U_0N))\). Take two
distinct primitive coefficient pairs with the same labels. Then
\(\ell_f\ell_g\mid f_1g_2-f_2g_1\). This coefficient determinant is a
nonzero integer, so its absolute value is at least \(\ell_f\ell_g\). As
\(U_0\le f_1,f_2<2U_0\), their ratios differ by at least
\(\ell_f\ell_g/(4U_0^2)\). Therefore
\[
 \#\{\text{coefficient pairs in one cell}\}
 \ll 1+\frac{U_0V_*}{N \ell_f\ell_g}.
\AutoTag{eq:old:7_17}
\]
The graph equation gives at most one positive primitive corner
pair from each coefficient pair. Fix an endpoint pattern.
Summing over the \(3N\) cells and the divisor-bounded labels gives
\[
 O_\eta\!\left(Y^\eta\left(N+
 \frac{U_0V_*}{\ell_f\ell_g}\right)\right).
\]
By the bound above for this pattern, the count is
\[
 O_\eta\!\left(Y^\eta\left(N+
 Y^{q-E_{\rm pat}+\Delta_{\rm exp}}\right)\right).
\]
The \(+1\) term, unit labels, and endpoint equalities are kept. The
eight chart-pattern classes have fixed cardinality. So their sum adds no
further power of \(Y\).
\end{enumerate}
The source corner coordinates are positive. The lattice census in
\eqref{eq:old:7_12} keeps both shortest-vector axes, which contribute
\(U_0+V_*\). In either counting orientation take \(c=c_1\) or
\(c=c_2\) respectively, as defined below. The four assigned factors of
\(f\) and the four complementary factors of \(g\) give
\[
 U_0\le f<2^4Y^c,\qquad V_*\le4g<2^6Y^{q-c}.
\]
These two terms are charged through the first two savings in
\eqref{eq:old:7_19}. The four graph patterns give its other four
savings. The grid and square-root terms are kept separately.

\subsection{The lower-corner linear program and lattice-point count}

The exact finite optimization below controls the source's choice of chart
or corner route. In each source box, let
\(x_0,\ldots,x_7\) be the lower endpoints of the base-\(Y\) logarithms
of the eight uniquely assigned corner factors, in the order
\[
 (a_A,a_B,b_C,b_D,c_A,c_C,d_B,d_D).
\]
Each factor is the full prime-power part of its residual supported on the
eligible corner set. Supported complementary parts are fixed by the complete
coefficient-corner labels and are charged through \eqref{eq:old:6_6}.
Write
\[
\begin{aligned}
x_{AB}&=x_0+x_1,& x_{AC}&=x_4+x_5,& q&=\sum_{i=0}^7x_i,\\
x_{ABCD}&=x_0+x_1+x_2+x_3,& c_1&=x_{AB}+x_{AC},\\
c_2&=x_0+x_1+x_6+x_7,&
P&=x_0+x_4,& Q&=x_3+x_7,\\
R_{\rm LP}&=x_1+x_6,& T&=x_2+x_5.
\end{aligned}
\AutoTag{eq:old:7_18}
\]
The source constraints include the canonical radical, shell, unit, chart,
and corner-capacity inequalities. Their limiting affine logarithmic forms
are listed in Appendix~\ref{app:arithmetic}. Their finite source collars
are recorded below. The source log coordinates
\(\mu,z_u,z_t\) also use their box lower endpoints. Thus the eleven
source coordinates are binned, and the auxiliary \(W\) is evaluated by
\eqref{eq:old:7_21} at those endpoints.
With \(\rho_Y\) as in \eqref{eq:old:7_8h}, every source-row inequality has
lower-corner slack at most
\(\rho_Y\). The 22 affine source rows in Appendix~\ref{app:arithmetic}
have coefficient \(\ell^1\)-norm at most \(8\). So replacing the
eleven source coordinates by reference values costs at most \(8d_2\).
In deriving these rows from \eqref{eq:old:7_8a}, the remaining finite terms
are logarithms of
\(n/Y\), \((m+n+1)/Y\), \((m+1)/m\), \((n+1)/n\), and
\(\kappa_Y=1+1/Y\). The first two have absolute logarithm at most
\(d_2\). The next two together are at most \(\ell_{\rm end}\). Each
source row uses at most three further \(d_2\) units, two
\(\ell_{\rm end}\) units, and two powers of \(\kappa_Y\). This gives the
stated \(\rho_Y\). The count is applied after these source-row comparisons.

For chart 1 put \((c,F,G)=(c_1,P,Q)\), and for chart 2 put
\((c,F,G)=(c_2,R_{\rm LP},T)\). Its six nominal savings are
\[
\begin{aligned}
 &\frac23-c,\quad \frac23-q+c,\\
 &\frac23-q+F+G-\Delta_{\rm exp},\quad
 \frac23-q+F+(q-c-G)-\Delta_{\rm exp},\\
 &\frac23-q+(c-F)+G-\Delta_{\rm exp},\quad
 \frac23-q+(c-F)+(q-c-G)-\Delta_{\rm exp}.
\end{aligned}
\AutoTag{eq:old:7_19}
\]
Let \(A_j\) be their minimum. The four assigned-corner savings are
\[
\begin{aligned}
 G_A&=\frac23-\left(\frac{\mu+z_u-c_1}{2}+x_1+x_5+z_u\right),\\
 G_D&=\frac23-\left(\frac{\mu+z_u-q+c_1}{2}+x_2+x_6+z_u\right),\\
 G_B&=\frac23-\left(\frac12-\frac{z_u}{2}+\frac q4-\frac{c_2}{2}
                   +x_0+x_7+z_t\right),\\
 G_C&=\frac23-\left(\frac12-\frac{z_u}{2}-\frac q4+\frac{c_2}{2}
                   +x_3+x_4+z_t\right).
\end{aligned}
\AutoTag{eq:old:7_20}
\]
At the lower corner set
\[
 W=\max(0,A_1,A_2,G_A,G_D,G_B,G_C).
\AutoTag{eq:old:7_21}
\]
The exact linear program has the twelve nonnegative coordinates
\((x_0,\ldots,x_7,W,z_u,z_t,\mu)\), 22 source rows, two chart rows
and four corner-route rows. The 72 branches come from the selected-chart
choice and the two six-way choices of minimizing chart scores. Dropping the
inequalities that enforce those choices enlarges each branch polytope.
Exact primal and dual elimination on these rational polytopes gives
\[
 W\ge L_{5W}-5\rho_Y,\qquad
 L_{5W}=\frac{1203113}{90000000}.
\AutoTag{eq:old:7_22}
\]
In every branch the counting-row dual mass is \(1\), the source-row dual
mass is at most \(5\), and the total dual mass is at most \(6\). The finite
certificate and its exact row matrix are given in Appendix
~\ref{app:arithmetic}.

Next we tie these nominal rows to counts on the same filtered box.
In chart 1 write \(\alpha=ac\), \(\beta=bd\). For \(G_A\), record
\((A,a_B,c_C,u)\), with \(a_Ac_A=f_J\) paid through
\eqref{eq:old:6_6}. The identity
\(u+\alpha A^2=\beta D^2\) chooses \(D\) by a square-divisor choice and
then fixes \(\beta\). For \(G_D\), record
\((D,b_C,d_B,u)\), with \(b_Dd_D=g_K\) paid in the same way. Then
\(\beta D^2-u=\alpha A^2\) chooses \(A\)
and then fixes \(\alpha\). In both cases split \(\alpha=ac,\beta=bd\)
and reconstruct
\[
 B=\frac{cA+bD}{u},\quad C=\frac{aA+dD}{u},\quad
 t=\frac{cd-ab}{u},\quad m=ABa,\quad n=ACc.
\]
In chart 2 write \(\gamma=ad\), \(\delta=bc\). For \(G_B\), record
\((B,a_A,d_D,t)\), with \(a_Bd_B\) paid through \eqref{eq:old:6_6}.
Then \(t+\gamma B^2=\delta C^2\) chooses \(C\)
and then fixes \(\delta\). For \(G_C\), record \((C,b_D,c_A,t)\), with
\(b_Cc_C\) paid through \eqref{eq:old:6_6}. Then
\(\delta C^2-t=\gamma B^2\) chooses \(B\) and fixes \(\gamma\).
Split \(\gamma=ad,\delta=bc\) and reconstruct
\[
 A=\frac{dB-bC}{t},\quad D=\frac{cC-aB}{t},\quad
 u=\frac{cd-ab}{t},\quad m=ABa,\quad n=ACc.
\]
Each square-divisor choice and supported-prime choice costs \(Y^\eta\)
at the fixed polynomial heights. The residual products must also be split
into their ordered corner factors. For example, splitting
\(\alpha=ac\) and \(\beta=bd\) costs at most
\(\tauDiv(\alpha)\tauDiv(\beta)\ll_\eta Y^\eta\). The same bound holds
for \(\gamma=ad\), \(\delta=bc\). Once these choices are fixed, the
inverse formulas give the other corners, and \(ut=L\) gives
the other shift. Thus the total inverse multiplicity is \(Y^\eta\).

Each encoding freely enumerates one corner, two individual assigned parts,
and the shift of that route. These four data have upper dyadic endpoints
costing \(4d_2\). The supported parts are determined by the full
coefficient-corner labels. They are not extra free variables. Supported
parts use the full prime-power convention of \eqref{eq:old:2_4}. Their
supported-prime choices are paid by the uniform estimate \eqref{eq:old:6_6}.
Each sum of two assigned-part coordinates is a lower bound for the
base-\(Y\) logarithm of their product.
To compare the corner with its LP reference value, put
\[
\nu_n=\log_Y n,\quad H=\log_Y(m+n+1),\quad
h_m=\log_Y(1+1/m),\quad h_n=\log_Y(1+1/n).
\]
Let subscript \(r\) denote the actual source values of \(\mu,z_u,q,c_1,c_2\).
Define the lower-corner reference exponents
\[
\begin{aligned}
\mathsf A&=(\mu+z_u-c_1)/2,&
\mathsf D&=(\mu+z_u-q+c_1)/2,\\
\mathsf B&=1/2-z_u/2+q/4-c_2/2,&
\mathsf C&=1/2-z_u/2-q/4+c_2/2.
\end{aligned}
\]
Take logarithms in \eqref{eq:old:7_8a}, and use \(Z^2M=m(m+1)n(n+1)\) and
\(ZL=(n-m)(m+n+1)\). This gives the exact formulas
\[
\begin{aligned}
\log_Y A&=(\mu_r+z_{u,r}-c_{1,r})/2+(\nu_n-H)/2,\\
\log_Y D&=(\mu_r+z_{u,r}-q_r+c_{1,r})/2
             +(\nu_n-H+h_m+h_n)/2,\\
\log_Y B&=1/2-z_{u,r}/2+q_r/4-c_{2,r}/2
             +(H-1)/2+(h_n-h_m)/4,\\
\log_Y C&=1/2-z_{u,r}/2-q_r/4+c_{2,r}/2
             +(H-1)/2+(h_m-h_n)/4.
\end{aligned}
\AutoTag{eq:old:7_23a}
\]
Each of the eleven independently binned source coordinates lies between
its lower endpoint and that endpoint plus \(d_2\). Apply this
coordinatewise to the linear forms in \eqref{eq:old:7_23a}, before
combining the corner coordinates into \(q,c_1,c_2\). This gives
\[
\begin{aligned}
\log_YA&\le\mathsf A+d_2,&
\log_YD&\le\mathsf D+3d_2+\ell_{\rm end}/2,\\
\log_YB&\le\mathsf B+3d_2/2+\lambda/2,&
\log_YC&\le\mathsf C+3d_2/2+\ell_C/2+\lambda/2.
\end{aligned}
\AutoTag{eq:old:7_23b}
\]
The exact correction in the \(C\) formula is
\((h_m-h_n)/4\le\ell_C/4\). The \(\ell_C/2\) shown is
conservative.
Here \(0<H-\nu_n\le d_2+\lambda\), \(H-1\le d_2+\lambda\),
\(h_m+h_n\le\ell_{\rm end}\), and
\(0\le h_m-h_n\le\ell_C\). The last bound follows from
\(m>Y^{\tau_0}/\sqrt3\), \(n>Y/2\), and \(m<n\). The coefficient
variation from the lower endpoints is at most \(d_2\) for \(\mathsf A\)
and \(3d_2\) for \(\mathsf D\). The \(3d_2\) bound for \(\mathsf D\) is
conservative. Cancellation between \(q\) and \(c_1\) reduces
the positive-coordinate variation in its lower-corner form to at most
\(d_2\). We keep the larger uniform allowance. In each of
\(\mathsf B,\mathsf C\),
the positive coefficients among \(q/4-c_2/2\) (or its negative) sum to
one across the original eight corner coordinates. So the upper variation is
at most \(d_2\), and the other coordinates decrease the
form. The other terms in \eqref{eq:old:7_23a} give the corrections shown.
Adding the \(4d_2\) data endpoints gives the
four route costs
\[
5d_2,\qquad 7d_2+\ell_{\rm end}/2,\qquad
11d_2/2+\lambda/2,\qquad
11d_2/2+\ell_C/2+\lambda/2.
\AutoTag{eq:old:7_23c}
\]
By \eqref{eq:old:7_3}, \(\ell_C\le\ell_{\rm end}\).
The two chart counts use the complete \eqref{eq:old:7_12} and
\eqref{eq:old:7_17}, with both axis terms included in the six chart
costs. They apply to the same filtered source box. Refine it at once
into the dyadic boxes of both counting orientations, keeping its original
selected-chart label. The fixed power of \(\log Y\) refinements is
absorbed in \(Y^\eta\). Their lower-corner loss is at most \(10d_2\).
This includes both axis terms and the four-part products in all graph
patterns. For \(Y\ge2\), every route cost is at most
\(\rho_Y\), by its definition and \(\ell_C\le\ell_{\rm end}\).
These are alternative bounds for the same source box.
Choose a route attaining \eqref{eq:old:7_21}. Its lattice-point count has
geometric saving at least \(W-\rho_Y\), before the arbitrary analytic
loss allocated in Section~\ref{sec:epsilon}. Then \eqref{eq:old:7_22}
gives \(L_{5W}-6\rho_Y\). The geometric,
interval, and square-root terms enter the completed source count
additively. The axis costs are already in the chart bound. The
number of terms is fixed. So the largest term exponent, that is the
minimum saving, controls their sum.

The interval and square-root terms in \eqref{eq:old:7_12} have savings
\[
\begin{aligned}
 s_{\rm int}&=\frac23-a_{\rm grid}\\
 &=\frac{235120131355662893}{9520800917171070000}
   \approx0.0246954,\\
 s_{\rm HD}&=\frac23-\frac{q_{\max}+a_{\rm grid}}{2}\\
 &=\frac{29390272487700347267}{4760400458585535000000}
   \approx0.00617391.
\end{aligned}\AutoTag{eq:old:7_24}
\]
For
\[
 c_N=\log_Y(1+Y^{-a_{\rm grid}}),\qquad E_q=2\lambda+\log_Y18,
\]
the interval term has saving at least \(s_{\rm int}-c_N\). For the
square-root term, \(M<Y^{q_{\max}+2\lambda}\) and
\(N<Y^{a_{\rm grid}+c_N}\) give
\[
 \sqrt{4MN}
 \le Y^{(q_{\max}+a_{\rm grid})/2+\lambda+c_N/2+d_2}.
\]
This is a conservative uniform upper bound for the finite high-degree
correction. It uses the largest allowed source product and the ceiling
bound for \(N\). It uses no correlation between those endpoints.
So the high-degree saving is at least
\(s_{\rm HD}-(E_q+c_N)/2-12d_2\), since
\[
 \frac{E_q+c_N}{2}+12d_2
 =\lambda+\frac{\log_Y18}{2}+\frac{c_N}{2}+12d_2
 \ge\lambda+\frac{c_N}{2}+d_2.
\]
Thus all endpoint terms are kept. Let
\(\eta>0\). Choose all finite determinant, factor, divisor,
Rankin, and logarithmic losses to total at most \(\eta/2\). Then
enlarge one threshold. This gives
\[
\sigma_{5W}=\min\!\left(L_{5W}-6\rho_Y,\,
 s_{\rm int}-c_N,\,
 s_{\rm HD}-(E_q+c_N)/2-12d_2\right).
\]
Thus
\[
 \#\{5W\text{ sources in one box}\}
 \ll_\eta Y^{\,2/3-\sigma_{5W}+\eta/2}.
\AutoTag{eq:old:7_25}
\]
The number of boxes is a fixed power of \(\log Y\). It is absorbed in
the remaining \(\eta/2\). The margins of the three limiting caps above
\(S_0\) are
\[
\begin{aligned}
 L_{5W}-S_0&=\begin{gathered}\frac{647471}{90000000}\\[-2pt]\text{\scriptsize (0.007194)}\end{gathered},\\
 s_{\rm int}-S_0&=\begin{gathered}\frac{88170305326616070517}{4760400458585535000000}\\[-2pt]\text{\scriptsize (0.018522)}\end{gathered},\\
 s_{\rm HD}-S_0&=\begin{gathered}\frac{128034121242821}{1190100114646383750000}\\[-2pt]\text{\scriptsize (1.076$\times10^{-7}$)}\end{gathered}>0.
\end{aligned}
\AutoTag{eq:old:7_26}
\]
The three limiting caps are at least \(S_0\) by \eqref{eq:old:7_26}. Their corrected
versions bound the additive contributions in \eqref{eq:old:7_25}, whose minimum controls
the sum exponent. For \(Y\ge4\), since \(\tau_0>1/2\),
\[
 (1+\sqrt3Y^{-\tau_0})(1+2/Y)(1+1/Y)
 <2\cdot\frac32\cdot\frac54=\frac{15}{4}<4.
\]
Thus \(\ell_{\rm end}+\lambda\le2d_2\), hence
\(\rho_Y\le15d_2\) and \(6\rho_Y\le90\log2/\log Y\).
The inequality \(\log(1+t)\le t\) for \(t\ge0\) and
\(a_{\rm grid}>1/2\) give, for \(Y\ge2\),
\[
 c_N\le\frac{Y^{-a_{\rm grid}}}{\log Y}<\frac1{\log Y},
 \qquad
 \lambda+\frac{c_N}{2}
 \le\frac{Y^{-1}+\tfrac12Y^{-1/2}}{\log Y}
 <\frac1{\log Y}.
\]
Therefore, for \(Y\ge4\),
\[
 \frac{E_q+c_N}{2}+12d_2
 \le\frac{\tfrac12\log18+12\log2+1}{\log Y}
 <\frac{90\log2}{\log Y}.
\]
Thus every correction in \eqref{eq:old:7_25} is at most
\(90\log2/\log Y\), and \eqref{eq:old:7_26} gives
\[
 Y^{S_0-\sigma_{5W}}\le2^{90}.
\]
The finite corrections therefore change only the implied constant.
After the box sum, \eqref{eq:old:7_25} implies the uniform estimate
\[
 \#5W\ll_\eta Y^{2/3-S_0+\eta}.
\AutoTag{eq:old:7_27}
\]
\begin{rowproposition}[\(5W\) row estimate]\label{prop:row-5w}
For every \(\eta>0\), the complete first-match remainder satisfies
\[
 \#5W\ll_\eta Y^{2/3-S_0+\eta}.
\]
\end{rowproposition}
\begin{proof}
The three additive source-box contributions are bounded in
\eqref{eq:old:7_25}. Their limiting caps and positive margins are
\eqref{eq:old:7_26}. This gives \eqref{eq:old:7_27}.
\end{proof}
The threshold in \eqref{eq:old:7_27} depends on \(\eta\) and the fixed
source parameters. It does not depend on any moving cell, shift, label,
factor, or resultant value. For each box and cell the degree set and
height exponents are fixed. For each selected factor the polynomial, its
endpoint coefficients and its resultant are fixed. They may vary between
boxes, cells and factors.

The exact global reserve is
\[
 g_0=S_0-\frac1{165}=\frac{18677}{165000000}>0.
\AutoTag{eq:old:7_28}
\]
It is used in Section~\ref{sec:epsilon} to choose one common loss and
one \(Y_{\rm final}\). The source-to-count argument above gives the
uniform estimate used there.
